\begin{table}[H]
\centering
\caption{Características de los hogares según condición de pobreza}
\scriptsize
\renewcommand{\arraystretch}{1.05}
\setlength{\tabcolsep}{4pt}
\begin{tabular}{p{5.2cm} p{3.3cm} p{3.3cm}}
\toprule
\textbf{Variable} & \textbf{No pobre (79.3\%)} & \textbf{Pobre (20.7\%)} \\
\textbf{}  & \textbf{Media (SD)} & \textbf{Media (SD)} \\
\midrule
N.º de cuartos & 3.50 (1.25) & 3.04 (1.13) \\
Cuartos para dormir & 1.99 (0.91) & 1.98 (0.84) \\
Personas en el hogar & 3.09 (1.65) & 4.15 (2.03) \\
Pago de arriendo (COP) & 494,501 (3,170,617) & 267,399 (2,288,991) \\
Edad promedio & 39.2 (16.6) & 30.6 (16.2) \\
Proporción de mujeres & 0.52 (0.29) & 0.55 (0.24) \\
Ocupación informal & 0.03 (0.15) & 0.03 (0.12) \\
Régimen subsidiado & 0.32 (0.39) & 0.58 (0.34) \\
Jóvenes “ninis” & 0.04 (0.12) & 0.09 (0.15) \\
Ocupados & 0.54 (0.32) & 0.31 (0.24) \\
Desocupados & 0.05 (0.14) & 0.08 (0.17) \\
Menores de 18 años & 0.18 (0.21) & 0.37 (0.24) \\
Menores de 5 años & 0.05 (0.12) & 0.11 (0.16) \\
Mayores de 65 años & 0.14 (0.29) & 0.10 (0.24) \\
Educación primaria & 0.04 (0.15) & 0.10 (0.20) \\
Educación secundaria & 0.02 (0.07) & 0.03 (0.09) \\
Educación superior & 0.23 (0.30) & 0.33 (0.29) \\
Empleo formal & 0.37 (0.42) & 0.56 (0.46) \\
\addlinespace
\textbf{Tipo de propiedad de la vivienda} & & \\
\quad Propia sin deuda & 0.41 & 0.28 \\
\quad Propia con deuda & 0.04 & 0.02 \\
\quad En arriendo & 0.37 & 0.43 \\
\quad Subarriendo & 0.15 & 0.16 \\
\quad Prestada / Usufructo & 0.03 & 0.11 \\
\quad Otra & 0.00 & 0.00 \\
\addlinespace
\textbf{Zona} & & \\
\quad Urbana & 0.91 & 0.86 \\
\quad Rural & 0.09 & 0.15 \\
\bottomrule
\multicolumn{3}{p{12.5cm}}{\footnotesize Las variables continuas presentan la media y desviación estándar (SD) entre paréntesis. Las variables categóricas muestran la proporción de observaciones dentro de cada grupo (n = 154.393) Los grupos se definen según la variable \texttt{pobre} (0 = No pobre, 1 = Pobre).}
\end{tabular}
\end{table}